\section{地下水・水理}

\subsection{平成24年度 問35}
\label{q:h24-reg-applied-35}

\noindent\textbf{原典：} sources/H24/17\_applied\_science.pdf、PDF物理10頁。

\textbf{公式正答：} 1（single）。

\par\textit{原典本文の校正と独自解説は未作成。分類のみの登録項目。}

\subsection{平成25年度 問29}
\label{q:h25-reg-applied-29}

\noindent\textbf{原典：} sources/H25/17\_applied\_science.pdf、PDF物理10頁。

\textbf{公式正答：} 1（single）。

\par\textit{原典本文の校正と独自解説は未作成。分類のみの登録項目。}

\subsection{令和元年度再試験 問29}
\label{q:r01-ret-applied-29}

\noindent\textbf{原典：} sources/R01\_retest/17\_applied\_science.pdf、PDF物理10頁。

\textbf{公式正答：} 5（single）。

\par\textit{原典本文の校正と独自解説は未作成。分類のみの登録項目。}

\subsection{令和7年度 問19}
\label{q:r07-reg-applied-19}

\noindent\textbf{原典：} sources/R07/17\_applied\_science.pdf、PDF物理7頁。

\textbf{公式正答：} 4（single）。

\par\textit{原典本文の校正と独自解説は未作成。分類のみの登録項目。}

\section{地下水の水質・水頭・含水量の定義を比較する。}

\subsection{平成26年度 問24}
\label{q:h26-reg-applied-24}

\noindent\textbf{原典：} sources/H26/17\_applied\_science.pdf、PDF物理10頁。

\textbf{公式正答：} 2（single）。

\paragraph{問題の要旨} 地下水の水質・水頭・含水量の定義を比較する。

\Needspace{4\baselineskip}\paragraph{確認した条件・記述（独自要約）}\mbox{}\par

\noindent\textbf{判断対象：} 地下水の水質・水頭・含水量の定義を比較する。\par

\noindent\textbf{原典確認：} 22原典頁画像で五候補・否定・化学式を照合。\par

\paragraph{問われる核心} 地下水の水質・水頭・含水量の定義を比較する。

\paragraph{必要知識} 地下水水質は起源の水だけでなく岩石との溶解・沈殿、酸化還元、生物・人為作用等を反映するので2。

\paragraph{考え方・正答を選ぶ理由（査読待ち初稿）} 地下水水質は起源の水だけでなく岩石との溶解・沈殿、酸化還元、生物・人為作用等を反映するので2。全水頭は位置と圧力の和、動水勾配は無次元。θは全体積を分母にし、飽和時θ=n、e=\allowbreak{}n/(1−n)。定義の名前だけでなく分母・次元を点検する。

\Needspace{4\baselineskip}\paragraph{選択肢の検討（査読待ち初稿）}\mbox{}\par

\noindent\textbf{1：選ばない候補} 帯水層に不飽和を含み被圧に自由水面。飽和し水を伝える帯水層と不飽和帯、自由水面の不圧を取り違え。\par

\noindent\textbf{2：選ぶ候補} 水質は降水/地表水/地層/環境に依存。起源と水岩反応・環境条件が水質を変える。\par

\noindent\textbf{3：選ばない候補} 全水頭は位置−圧力、勾配を拡散率。低速地下水のh=\allowbreak{}z+p/(ρg)は和で、差/距離は動水勾配。\par

\noindent\textbf{4：選ばない候補} 体積含水率は水/間隙、飽和で間隙比。θ=Vw/V全体、飽和では間隙率n。間隙比eではない。\par

\noindent\textbf{5：選ばない候補} 基底流は低水期除外、都市で直接流出減。基底流は低水期にも地下水等が支える流量。都市流出の断定も逆。\par

\paragraph{再利用する考え方} 原子番号と周期・族、物性と反応性、酸強度と危険性、生成と除去、定義の分母を別々に確かめる。

\paragraph{共通理論} \hyperref[theory:hydro:definitions]{第\ref*{theory:hydro:definitions}節}

\par 参考：\cite{inorganic_waterchem,inorganic_hydro_ref}。

\paragraph{関連問題} \hyperref[q:h27-reg-applied-23]{H27-REG-APPLIED-23}

\par 難易度 2（校準初稿）、優先度未評価。

\subsection{平成27年度 問23}
\label{q:h27-reg-applied-23}

\noindent\textbf{原典：} sources/H27/17\_applied\_science.pdf、PDF物理9頁。

\textbf{公式正答：} 3（single）。

\par\textbf{検討未確定：} 大気/河川平均滞留時間は資料・期間・推計の定義に依存。原典8日とUSGS約10日の違いを保持。

\paragraph{問題の要旨} 含水比の分母と陸水の移動過程を比較する。

\Needspace{4\baselineskip}\paragraph{確認した条件・記述（独自要約）}\mbox{}\par

\noindent\textbf{判断対象：} 含水比の分母と陸水の移動過程を比較する。\par

\noindent\textbf{原典確認：} 22原典頁画像で五候補・否定・化学式を照合。\par

\paragraph{問われる核心} 含水比の分母と陸水の移動過程を比較する。

\paragraph{必要知識} 通常の土質含水比は水質量/乾燥固体質量で、湿潤土全体で割った3は別の質量分率。

\paragraph{考え方・正答を選ぶ理由（査読待ち初稿）} 通常の土質含水比は水質量/乾燥固体質量で、湿潤土全体で割った3は別の質量分率。同じ試料でも値が異なる。失水/得水河川は交換の向き、Horton流は浸透能力を越える雨、蒸発散は気化の合算。大気滞留約8日の値を全推計の厳密定数とせず、比較関係と数値の留保を分ける。

\Needspace{4\baselineskip}\paragraph{選択肢の検討（査読待ち初稿）}\mbox{}\par

\noindent\textbf{1：選ばない候補} 大気平均滞留8日で河川より短い。約8日は推計条件依存。USGSには大気約10日・河川16–26日例もある。\par

\noindent\textbf{2：選ばない候補} 失水河川は地下へ浸透。河川から帯水層へ水が失われる区間。\par

\noindent\textbf{3：選ぶ候補} 含水比は水/湿潤土全体。通常土質含水比w=\allowbreak{}Mw/M乾燥であり湿潤全体ではない。\par

\noindent\textbf{4：選ばない候補} 浸透能より降雨強度大ならHorton流。浸透能力を超え余剰が地表へ流れる。\par

\noindent\textbf{5：選ばない候補} 蒸散と蒸発を合わせ蒸発散。植物経由と非植物面の気化を合算。\par

\paragraph{再利用する考え方} 原子番号と周期・族、物性と反応性、酸強度と危険性、生成と除去、定義の分母を別々に確かめる。

\paragraph{共通理論} \hyperref[theory:hydro:definitions]{第\ref*{theory:hydro:definitions}節}

\par 参考：\cite{inorganic_hydro_ref,inorganic_waterbudget}。

\paragraph{関連問題} \hyperref[q:h26-reg-applied-24]{H26-REG-APPLIED-24}

\par 難易度 2（校準初稿）、優先度未評価。
